% Prospective journal protocol; replace planning language after experiments.
\begin{journaladdition}
\subsection{Experimental Design and Model Selection}
\label{sec:experiments}
\label{sec:model_selection}
All experiments below are planned. Settings will be selected on validation patients and frozen before evaluation on test patients.
The split, endpoints, and analysis rules will be fixed before running the core experiments.
\todo{Define the validation selection criterion, sweep ranges, compute budgets, and repeat/seed policy. Do not reuse the previous surrogate-set ranking as journal test evidence.}

\subsubsection{PET and rigidity ablations.}
\begin{itemize}
  \item Train the $2\times2$ conditions: neither PET preservation nor rigidity, PET only, rigidity only, and both. Retain the same registration-accuracy and general regularization terms.
  \item Within the full objective, remove each of $\mathcal{L}_{\mathrm{MTV}}$, $\mathcal{L}_{\mathrm{MTV\text{-}J}}$, $\mathcal{L}_{\mathrm{jac}}$, and $\mathcal{L}_{\mathrm{TLG}}$ individually.
  \item Prioritize matched training for the central claims. Specify each condition's training and IO objective separately; identify any IO-only ablation as such.
  \item Report paired effects of PET, rigidity, and their combination on alignment, lesion preservation, and deformation quality; include the PET--rigidity interaction if estimable under the chosen analysis.
\end{itemize}

\subsubsection{Comparison of rigidity penalties.}
Compare no rigidity, Staring with all three components~\cite{ref_rigidity_3_staring}, Staring with mask erosion, and the existing Kabsch rigid-fit penalty~\cite{ref_kabsch}.
Tune or sweep each method's weight separately; equal numerical weights need not imply comparable constraint strength.
\begin{itemize}
  \item Document the three Staring components, their scaling, derivative stencils, boundary handling, and coordinate units. Specify the erosion rule and report bones/voxels retained.
  \item Use alignment versus physical within-bone distance distortion as the main comparison, with folding alongside. Report bone rigidity error (BRE) only as a secondary endpoint because it closely matches the Kabsch objective.
  \item Measure direct stencil leakage separately from the effect on surrounding tissue. Kabsch avoids direct stencil leakage, but shared deformation parameters and regularization can still influence soft tissue.
  \item Scope conclusions to the tested resolutions, implementations, and weight ranges.
\end{itemize}

\subsubsection{Accuracy--preservation trade-off.}
Run a small PET-weight sweep including zero and the current setting.
\todo{Specify whether one multiplier scales the whole PET group, which relative term weights remain fixed, and whether each setting is retrained or changes IO only.}
Plot alignment against lesion preservation and report deformation quality for every setting.
Select operating points on validation patients and evaluate frozen settings on the test partition, without choosing a preferred setting from test results.

\subsubsection{Baselines.}
\label{sec:baselines}
Re-evaluate the initial alignment, affine registration, NiftyReg~\cite{MODAT2010,modat2014}, and ConvexAdam~\cite{ref_convexadam} on the new test partition under the same evaluation definitions.
Report the full method with and without IO.
\todo{Freeze baseline configurations, validation tuning allowances, and runtime measurement. The previous container and surrogate-validation configurations are not journal comparison groups.}

\subsection{Evaluation Endpoints and Statistical Analysis}
\label{sec:evaluation}
\begin{itemize}
  \item \textbf{Alignment:} per-bone and soft-tissue organ Dice and surface distance; near-bone landmarks only if independently available. Define aggregation over structures and handling of absent or truncated anatomy.
  \item \textbf{PET preservation:} aggregate and per-lesion MTV/TLG change from the original moving scan to that same scan after warping. Report signed, absolute, and upper-tail errors, stratified by bone association and lesion size.
  \item \textbf{Lesion measurement:} define lesions on the source scan and track them through the transformation. Predefine size bins, bone-association rules, minimum denominators, empty-mask handling, field-of-view exclusions, SUV units, and interpolation. Distinguish resampling effects from deformation effects using a specified control.
  \item \textbf{Bone shape:} measure distortion of sampled within-bone point-pair distances in physical coordinates, covering short and long distances. Predefine sampling, distance bins, and normalization; report BRE secondarily.
  \item \textbf{Deformation quality:} report NDV/folding inside bones, in an outside boundary band, and across the body. Specify masks, band width in physical units, denominators, and consistent units.
  \item \textbf{Boundary mechanism:} quantify the fraction of stencil samples outside bone and the fraction of bones/voxels retained after erosion. State how bones that disappear after erosion are handled.
  \item \textbf{Uncertainty:} report paired method effects with patient-level confidence intervals, accounting for multiple pairs and lesions per patient. A candidate is patient-cluster bootstrap resampling that retains all within-patient observations and method pairing.
\end{itemize}
\todo{Freeze primary endpoints and contrasts, the upper-tail statistic, patient weighting, uncertainty procedure, and multiplicity handling before test evaluation. Report evaluable patient/pair/lesion counts and exclusions.}

\subsection{Optional Supporting Analyses}
\label{sec:optional}
\begin{itemize}
  \item \textbf{Gradient analysis:} one loss-pair cosine matrix, selected IO trajectories, and representative maps. Distinguish displacement-space maps from gradients in optimized velocity coordinates; mask near-zero gradients and report weighted magnitudes. Local cosine similarity is not evidence of a final performance trade-off.
  \item \textbf{External cohort:} include only if feasible. Combine automatic measurements with blinded, randomized clinician review; add sparse landmarks if available. Separate IO labels from independent evaluation. Distinguish resampling from deformation effects.
\end{itemize}
These analyses are optional; independent evaluation and the core experiments take priority.
\end{journaladdition}
